\documentclass[../../../thesis.tex]{subfiles}
\begin{document}
\subsection{Grafy}
Grafmi budeme nazývať zobrazenie vedeckých dát
najčastejšie vo forme dvojrozmerného grafu
závislosti dvoch alebo viacerých veličín.
Príklad takéhoto objektu je na obrázku \ref{fig:Graph1}.
Grafická reprezentácia vedeckých dát musí byť
v~prvom rade čitateľná, zreteľná a~jednoznačná.
Tomu treba prispôbiť všetky zásady pri tvorbe grafov.

\begin{figure}[h!]
  \centering
  \includegraphics[scale=0.358]{img/Graph1}
  \caption{Ukážka grafu vytvoreného v externom progрame a vloženého ako PDF súbor.
  Použité písmo je Arial s veľkosťou približne 10\,pt. Plné krúžky sú body merania a~prerušovaná čiara je kvadratický fit závislosti $s = at^2/2$, pričom $a = (2{,}00 \pm0{,}01)\,\mathrm{m\,s^{-2}}$.}
  \label{fig:Graph1}
\end{figure}

\subsubsection{Formát súboru}
Vektorové formáty SVG alebo PDF sú ideálna voľba pri exporte
z~grafických progрamov, napr. z~Excelu alebo Originu.
Ak takúto možnosť nemáme, treba grafy z externého softvéru
exportovať do bitmapového formátu, najlepšie PNG.
Stratový formát JPEG nie je na čiarovú grafiku vhodný.
Rozlíšenie bitmapového súboru by malo byť minimálne 600 dpi,
aby boli čiary ostré.
Znamená to, že ak predpokladáme veľkosť obrázku
$10\,\mathrm{cm} \times 7{,}5\,\mathrm{cm}$,
musí mať aspoň 2\,363\,px $\times$ 1\,772\,px (pixelov).

\subsubsection{Písmo a hrúbka čiar}
Písmo v grafe nemusí byť nevyhnutne Computer Modern.
V~obrázkoch a~schémach sa často používa
tzv. bezserifové alebo groteskové písmo ako napr. Arial,
ktoré je lepšie čitateľné.
Veľkosť písma v~obrázkoch by nemala byť menšia než
10\,pt,
čo je o~dva stupne menej ako základná veľkosť písma
v~dokumente.

Pozornosť treba venovať aj dostatočnej hrúbke čiar osí
a~grafického znázornenia dát,
aby boli viditeľné aj po vytlačení na bežnej tlačiarni.

\subsubsection{Prvky grafu}
Formálne prvky grafu sú osi s~dielikmi a~číslami,
názvy osí s~uvedením veličín, násobkov a~jednotiek,
mriežka a~legenda.
Medzi obsahové prvky zaraďujeme znázornené hodnoty vo forme
bodov alebo čiar.
Graf môže obsahovať aj názov grafu
a~doplňujúce texty,
prípadne ďalšie grafické prvky na zvýraznenie niektorých bodov,
oblastí a~podobne.

Bežný graf pozostáva zväčša z~dvoch navzájom kolmých číselných
osí – z~ľavej zvislej a~spodnej vodorovnej,
ktoré sa pretínajú v~ľavom dolnom rohu.
Na spodnej osi sa nachádzajú hodnoty nezávislej veličiny,
ľavá zvislá os obsahuje hodnoty závislej veličiny.
Rozsahy osí volíme tak,
aby korešpondovali s~intervalmi zobrazovaných hodnôt,
prípadne aby znázorňovali javy,
ktoré majú byť z grafu zrejmé.
Osi sa môžu pretínať aj v inom než nulovom bode.

\subsubsection{Označenie osí}
Osi musia byť riadne označené názvom alebo značkou veličiny,
jej jednotkou a~násobkom.
Nedodržanie tohto pravidla sa považuje za závažný nedostatok
a~autor musí mať na takýto krok obhájiteľný dôvod.
Jednotku spolu s~násobkom uzatvárame kvôli jednoznačnosti
do okrúhlych zátvoriek.
Hranaté zátvorky sa v knižnej tlači na tento účel nepoužívajú.

Os musí byť jasne rozdelená dielikmi,
ktoré sú kolmé na os a~predstavujú okrúhle hodnoty
zobrazovanej veličiny.
V blízkosti hlavných dielikov sa nachádzajú čísla prislúchajúce
hodnote dieliku.
Táto hodnota sa potom násobí s údajom v~zátvorke
v~opise osi a~spolu tvoria hodnoty zobrazenej fyzikálnej
veličiny aj s~jednotkou.

\subsubsection{Viacero grafov v jednom obrázku}
Priebehy dvoch a~viac nezávislých veličín môžeme nakresliť
do spoločných osí alebo použijeme pravú nezávislú zvislú os.
V~špeciálnych prípadoch môžeme využiť aj hornú vodorovnú os.
Ak chceme v~jednom obrázku zobraziť viacero grafov,
musí byť príslušnosť jednotlivých bodov a~čiar k~osiam jasná
z~legendy.
Legendu možno zahrnúť aj do textu pod obrázkom.

Graf znázorňujúci experimentálne hodnoty fyzikálnych veličín
zvykne byť uzavretý zhora aj sprava tak,
ako na obrázku \ref{fig:Graph1}.
Dve prekrížené otvorené osi sa používajú zväčša
v~prípade teoretického nákresu matematickej funkcie $y = f(x)$.
\end{document}